\section{Recovering a latent invariant}
\label{sec:setup}

\paragraph{Model and data.} We analyze the published DreamerV3 architecture
\citep{hafner2023dreamerv3}: 13.5M parameters, categorical $32\times32$ stochastic latents, a
convolutional encoder and decoder, KL balancing, and unimix. We train only the world model, with no
actor or critic.

Data come from Gymnasium's \texttt{Pendulum-v1} \citep{towers2024gymnasium}, rendered per step and
block-averaged to $64\times64$, with the simulator state set directly
($\theta\sim U[-1.8,1.8]$, $\dot\theta\sim U[-2.2,2.2]$) and any trajectory reaching the
$|\dot\theta|=8$ clip rejected, since clipping breaks conservation. Actions are zero unless stated,
and a single frame does not reveal velocity, so the model must infer $\dot\theta$ from the sequence.
We train on 204 trajectories of 120 frames for about 6{,}500 gradient steps, reserving 52 as a
post-training \emph{analysis set}. Three seeds pass checks fixed before analysis: KL above 1 nat,
one-step decoding $4\times$ better than the dataset mean, and finite rollouts;
Appendix~\ref{sec:details} gives the rest. Our claims concern a DreamerV3
world model trained on pendulum video, not DreamerV3 as a full reinforcement-learning agent.

\paragraph{Latent coordinates.} After training we freeze the model. Let $h$ be the deterministic
recurrent state on the analysis set and $\bar h$ its mean. We take the top 12 principal directions of
$h$, drop any without support in the data, and let $P$ be the resulting orthonormal map. All
extraction and perturbation operate on $z = P^\top(h-\bar h)$. Let $F(z)$ be the projection through
$P$ of the model's one-step displacement $T(h)-h$; using the model's own displacement avoids fitting
a surrogate vector field and then analysing that instead.

\paragraph{The candidate family.} We search degree-4 polynomials in $z$: with $\varphi(z)$ the
monomials up to degree 4 in the 12 latent coordinates, a candidate is $C(z) = a^\top\varphi(z)$ over
1819 monomials, excluding the trivially conserved constant. Energy is quadratic in $\dot\theta$ and
needs polynomial terms to approximate $\cos\theta$, and the latent is an unknown nonlinear encoding
of $(\theta,\dot\theta)$ rather than those coordinates, so the family must express energy through
that distortion while staying small enough to fit.

\paragraph{The invariance criterion.} An invariant is constant along a trajectory and differs between
trajectories, so we score a candidate by
\begin{equation}
\text{ratio}(C) \;=\;
\frac{\text{mean within-trajectory variance of } C}{\text{total variance of } C},
\label{eq:ratio}
\end{equation}
which is $0$ for a perfectly conserved scalar and near $1$ for one wandering as much within a
trajectory as across the dataset. The denominator keeps it non-trivial: without it $C=0$ would score
perfectly while distinguishing nothing. Both variances are quadratic forms in $a$, so with $W$ the
mean within-trajectory covariance and $T$ the total covariance this is a Rayleigh quotient
$a^\top W a / a^\top T a$, minimised by the smallest generalized eigenvector of $W a = \lambda T a$.
One eigendecomposition returns the whole family, ranked.

\paragraph{Selecting among conserved candidates.} Conservation alone does not pick out a unique
scalar: if $C$ is conserved so is any function of it, so $E$, $E^2$ and mixtures all sit near the top
and the leading eigenvector is generally a combination. We therefore fit $C$ jointly with an
antisymmetric operator $B$ such that $F \approx B\nabla C$, $B^\top = -B$, mirroring the Hamiltonian
relation between a conserved quantity and the flow it generates \citep{arnold1989mathematical}. Since
$\nabla C^\top B\nabla C = 0$ for antisymmetric $B$, such a scalar is constant along the flow by
construction, and the criterion separates candidates the ratio cannot --- $E^2$ is exactly as
conserved as $E$ but does not pair with the same $B$. The fit is bilinear, so we alternate
least-squares solves, searching within the top eight eigenvectors so the recovered $C$ has eight
effective degrees of freedom rather than 1819 --- which matters on 52 trajectories, where a free fit
could drive the in-sample ratio to zero by overfitting. Appendix~\ref{sec:details} isolates its
contribution: conservation drives most of the recovery and intervention effect, while flow alignment
mainly reduces variation across seeds.

\paragraph{Reference energy.} We compare $C$ against textbook pendulum energy
$E = \tfrac{1}{6}\dot\theta^2 + 5\cos\theta$. Gymnasium integrates semi-implicitly, so it does not
conserve $E$ exactly; Section~\ref{sec:mechanism} shows what it conserves instead, and that this is
the crux rather than a caveat. The search sees only latent states and the model's own one-step
transition --- never $\theta$, $\dot\theta$ or $E$. Ground truth enters afterward, when scoring $C$
against $E$. The extraction dimension was selected on three development models and fixed before
evaluating the three reported here.
